\section{Workforce consequences: fewer FTE for the same governed output}
\label{sec:workforce}

The purpose of automating a gate is not a faster document. It is that the same governed output
requires fewer people. This section states that consequence without euphemism and marks exactly what
it depends on.

\subsection{From vanishing human effort to vanishing posts}

\textbf{Composition.} If each component satisfies its contract on every valid history it can receive,
outputs satisfy successor input requirements, and the controller and recovery are also automated, the
route executes without a human step. This is ordinary composition of compatible procedures. With at
most $M$ visits and a valid-prefix conditional failure bound $\epsilon$, summing first-failure events
gives a route bound $M\epsilon$, without an independence assumption. An average benchmark result
supplies neither that conditional bound nor a guarantee on arbitrary new inputs.

\textbf{Labor.} For a finite registry of packages, let $\nu_i(t)$ be visits, $\ell_i(t)$ mean human
hours per visit, and $B(t)$ nonduplicated recurring support. Then
\begin{equation}
L(t)=\sum_i \nu_i(t)\ell_i(t)+B(t),\qquad N(t)=L(t)/K. \label{eq:laborpath}
\end{equation}
At fixed output demand with bounded visits, $\ell_i(t)\to 0$ and $B(t)\to B_\infty$ imply
$N(t)\to B_\infty/K$. If support is also automated, $B_\infty=0$ and the requirement tends to zero. The
arithmetic consequence is direct: the same completed work requires fewer workers' hours, and complete
closure removes the execution labor requirement. A declining share of escalated cases is not enough:
$q(t)=1/t$ with exception hours $h_E(t)=t$ leaves expected exception labor at one hour per case. Exact
zero staffed execution requires removal of every necessary human task, not rounding a small FTE to
zero; a positive ten-minute obligation can still require an available person.

\textbf{Posts.} With interchangeable workers, useful capacity $uK$ per person at target utilization
$0<u<1$, positive wage cost, and fixed nonlabor costs, cost-minimizing headcount is
\begin{equation}
n^{*}(t)=\left\lceil \frac{L(t)}{uK} \right\rceil . \label{eq:headcount}
\end{equation}
For $n_0^{*}=\lceil L_0/(uK)\rceil$, a reduction of at least one post occurs exactly when
$L(t)\le (n_0^{*}-1)uK$. At complete closure $L=0$ and $n^{*}=0$. Reduced staffing is thus a direct
implication once productive necessity and the staffing objective are specified. Service levels and
heterogeneous skills add the requirements of Section~\ref{sec:staffing}.

An employer can retain people, shorten working time, or assign new work. Those choices do not restore
demand for the displaced gate execution. If the organization no longer needs ten people to operate the
same Procurement package, employing those ten elsewhere is a new allocation of labor, not evidence
that the old gate still needs them. The direct conclusion is a reduction of the workers required by
the replaced workflows; actual workforce reductions follow when staffing adjusts instead of the
released capacity being absorbed by new demand. In the author's experience, that adjustment is a stated
objective of some governance automation programs.

\subsection{The trajectory and its support floor}

Appendix~\ref{app:trajectory} allocates the 140 FTE of Table~\ref{tab:populations} across the three
routes and computes illustrative substitution configurations. With 30\%, 75\%, and 25\% of baseline
route labor remaining on W1, W2, and W3 and five FTE of shared support, the requirement is 60.56 FTE;
with 5\%, 20\%, and 3\%, it is 16.35; with direct gate execution closed, the five support FTE remain,
a reduction of 96.4\%; with support closed, it is zero. These are chosen fractions, not an estimated
trajectory: they show what the account returns, not when.

A fixed support requirement prevents a zero-worker endpoint even after every direct case activity is
automated. Law~5 predicts substitution of this remaining work as well; the model does not make the
floor vanish by renaming its employees supervisors. Conversely, a small supervisor team is already a
large structural reduction from a large case-processing population. Section~\ref{sec:example} shows the
other side of the same account: before the endpoint, badly governed automation can \emph{raise} the
requirement above 140.

\subsection{Employment beyond the DGF}

\citet{acemoglu2019} distinguish displacement from new-task reinstatement and productivity effects.
Those mechanisms can offset losses in particular activities. If output demand scales by $g(t)$ while
labor per output is $h(t)$, total variable labor follows $g(t)h(t)$: halving labor intensity while
doubling demand leaves it unchanged. New work on agent governance can add to $B(t)$ or create a
genuinely new output class. Including it is not softening displacement; it is measuring whether the
displaced labor is actually replaced by another demand.

A permanently positive support floor defeats the
stronger zero-labor closure even if the original review workforce is much smaller. The paper does not
equate either outcome with an automatic economy-wide reduction in employment; it does claim that the
governance specialists of the DGF are among the first white-collar populations for whom the question
stops being hypothetical.
